% ============================================================
% CHAPITRE 26 — Photomontage publicitaire avancé
% ============================================================

\chapter{Photomontage publicitaire avancé}

\noindent
\textit{Nous voici au dernier chapitre de la quatrième partie, consacrée au design publicitaire et à la communication visuelle. Au fil des chapitres précédents, nous avons appris à conduire un projet de design graphique, à créer des affiches et des flyers, à concevoir des visuels pour les réseaux sociaux, à construire un système d'identité visuelle, à élaborer une charte graphique. Toutes ces compétences convergent vers l'un des exercices les plus aboutis et les plus exigeants du design graphique : le \emph{photomontage publicitaire}.}

\vspace{1em}

\noindent
Le \terme{photomontage publicitaire} est l'art de créer une image qui n'existe pas, en assemblant plusieurs photographies, illustrations, textures, pour construire un univers visuel cohérent, séduisant, mémorable. Il est utilisé dans la publicité, le cinéma, l'édition, le packaging, le web. Il permet de mettre en scène un produit dans un environnement imaginaire, de créer une situation surréaliste, de raconter une histoire, de susciter une émotion. C'est l'un des domaines les plus créatifs et les plus techniques du design graphique.

\vspace{1em}

\noindent
Nous avons déjà abordé le photomontage dans la deuxième partie, à travers les chapitres sur Photoshop. Nous avons appris à détourer, à harmoniser les dimensions et les perspectives, à adapter la lumière et les couleurs, à créer des ombres, à ajouter des effets de profondeur. Mais le photomontage publicitaire va plus loin : il exige une \emph{maîtrise avancée} de tous ces outils, une \emph{compréhension fine} du concept publicitaire, une \emph{rigueur} dans l'intégration du produit, une \emph{créativité} dans la mise en scène. C'est cette dimension avancée que nous allons explorer dans ce chapitre.

\vspace{1em}

\noindent
Ce chapitre a pour objectif de te faire maîtriser le photomontage publicitaire avancé. Nous commencerons par l'analyse du concept publicitaire. Nous apprendrons ensuite à rechercher et sélectionner les photographies, à construire la composition, à réaliser un détourage avancé, à harmoniser la lumière, à gérer les ombres et les reflets, à intégrer le produit, à traiter la colorimétrie, à intégrer la typographie, à finaliser et à contrôler la qualité. À la fin de ce chapitre, tu sauras créer des visuels publicitaires spectaculaires, crédibles et professionnels.

% ============================================================
\section{Analyse du concept publicitaire}
% ============================================================

\subsection{Qu'est-ce qu'un concept publicitaire ?}

\begin{quote}
\textbf{Un concept publicitaire est l'idée centrale d'une campagne : la manière dont le produit, le message et le public sont mis en relation pour produire un effet. Il est l'âme du photomontage.}
\end{quote}

Le concept publicitaire est la \terme{colonne vertébrale} du photomontage. Sans concept, le photomontage n'est qu'une accumulation d'images ; avec concept, il devient une histoire, une émotion, un message.

\subsection{Les composantes d'un concept}

\subsubsection{L'idée}

L'\terme{idée} est le cœur du concept. Elle doit être :

\begin{itemize}
    \item \terme{Forte} : elle doit marquer.
    \item \terme{Simple} : elle doit être comprise immédiatement.
    \item \terme{Pertinente} : elle doit correspondre au produit, au public, à la marque.
    \item \terme{Originale} : elle doit se distinguer.
    \item \terme{Crédible} : elle doit être cohérente avec la réalité du produit.
\end{itemize}

\subsubsection{Le message}

Le \terme{message} est ce que le concept veut dire. Il doit être :

\begin{itemize}
    \item \terme{Clair} : compris immédiatement.
    \item \terme{Unique} : un seul message.
    \item \terme{Percutant} : mémorable.
    \item \terme{Cohérent} : avec la marque.
\end{itemize}

\subsubsection{Le bénéfice}

Le \terme{bénéfice} est ce que le public retire du produit. Il peut être :

\begin{itemize}
    \item \terme{Fonctionnel} : ce que le produit fait.
    \item \terme{Émotionnel} : ce que le produit fait ressentir.
    \item \terme{Symbolique} : ce que le produit représente.
\end{itemize}

\subsubsection{Le ton}

Le \terme{ton} est la manière dont le message est transmis. Il peut être :

\begin{itemize}
    \item \terme{Sérieux} : institutionnel, professionnel.
    \item \terme{Ludique} : humoristique, décalé.
    \item \terme{Élégant} : raffiné, sophistiqué.
    \item \terme{Dramatique} : intense, émouvant.
    \item \terme{Inspirant} : motivant, positif.
\end{itemize}

\subsubsection{L'univers visuel}

L'\terme{univers visuel} est l'ensemble des éléments graphiques qui traduisent le concept : couleurs, lumières, textures, ambiances, compositions.

\subsection{Analyser le concept}

\subsubsection{Étape 1 : Lire le brief}

\begin{itemize}
    \item \terme{Comprendre} le produit.
    \item \terme{Comprendre} le public.
    \item \terme{Comprendre} le message.
    \item \terme{Comprendre} le bénéfice.
    \item \terme{Comprendre} le ton.
\end{itemize}

\subsubsection{Étape 2 : Identifier l'idée}

\begin{itemize}
    \item \terme{Quelle} est l'idée centrale ?
    \item \terme{Comment} est-elle exprimée ?
    \item \terme{Pourquoi} est-elle pertinente ?
\end{itemize}

\subsubsection{Étape 3 : Identifier les éléments visuels}

\begin{itemize}
    \item \terme{Quel} est le produit ?
    \item \terme{Quel} est l'environnement ?
    \item \terme{Quels} sont les éléments secondaires ?
    \item \terme{Quel} est l'éclairage ?
    \item \terme{Quelles} sont les couleurs ?
\end{itemize}

\subsubsection{Étape 4 : Traduire en images}

\begin{itemize}
    \item \terme{Traduire} l'idée en images.
    \item \terme{Traduire} le message en visuel.
    \item \terme{Traduire} le ton en ambiance.
\end{itemize}

\subsubsection{Étape 5 : Vérifier la cohérence}

\begin{itemize}
    \item \terme{Vérifier} que le concept est cohérent avec la marque.
    \item \terme{Vérifier} que le concept est adapté au public.
    \item \terme{Vérifier} que le concept est réalisable.
\end{itemize}

\subsection{Les erreurs courantes}

\subsubsection{Pas de concept}

Un photomontage sans concept est une simple accumulation d'images.

\subsubsection{Concept faible}

Un concept banal, générique, sans originalité.

\subsubsection{Concept incohérent}

Un concept qui ne correspond pas au produit, au public, à la marque.

\subsubsection{Concept irréalisable}

Un concept impossible à réaliser techniquement, ou trop coûteux.

\subsection{Les bonnes pratiques}

\begin{itemize}
    \item \terme{Lire} attentivement le brief.
    \item \terme{Identifier} l'idée.
    \item \terme{Identifier} les éléments visuels.
    \item \terme{Traduire} en images.
    \item \terme{Vérifier} la cohérence.
    \item \terme{Vérifier} la faisabilité.
    \item \terme{Documenter} le concept.
    \item \terme{Présenter} au client.
\end{itemize}

\subsection*{Conclusion de la section}

Un \terme{concept publicitaire} est l'idée centrale d'une campagne. Il comprend une idée, un message, un bénéfice, un ton, un univers visuel. L'analyse suit plusieurs étapes : lire le brief, identifier l'idée, identifier les éléments visuels, traduire en images, vérifier la cohérence. Il est essentiel d'avoir un concept fort, cohérent, réalisable.

Maintenant que nous savons analyser le concept, nous allons, dans la section suivante, apprendre à \terme{rechercher et sélectionner des photographies}.

% ============================================================
\section{Recherche et sélection des photographies}
% ============================================================

\subsection{Pourquoi la recherche de photographies ?}

La \terme{recherche de photographies} est essentielle pour :

\begin{itemize}
    \item \terme{Constituer} l'univers visuel.
    \item \terme{Trouver} les images adaptées au concept.
    \item \terme{Assurer} la cohérence des images.
    \item \terme{Respecter} les droits d'auteur.
    \item \terme{Optimiser} la qualité.
    \item \terme{Gagner} du temps.
\end{itemize}

\subsection{Les sources de photographies}

\subsubsection{Les banques d'images}

\begin{itemize}
    \item \terme{Unsplash, Pexels, Pixabay} : images gratuites, libres de droits.
    \item \terme{Adobe Stock, Shutterstock, Getty Images} : images payantes, de haute qualité.
    \item \terme{Stocks spécialisés} : images de produits, de personnes, de lieux.
\end{itemize}

\subsubsection{Les photographies personnelles}

\begin{itemize}
    \item \terme{Photos} prises par vous-même.
    \item \terme{Shootings} réalisés pour le projet.
    \item \terme{Images} fournies par le client.
\end{itemize}

\subsubsection{Les images générées}

\begin{itemize}
    \item \terme{Intelligence artificielle} (Midjourney, DALL·E, Adobe Firefly) : images générées.
    \item \terme{Images} de synthèse.
\end{itemize}

\subsection{Les critères de sélection}

\subsubsection{La compatibilité}

Les photographies doivent être \terme{compatibles} entre elles :

\begin{itemize}
    \item \terme{Résolution} comparable.
    \item \terme{Lumière} cohérente (direction, intensité, couleur, dureté).
    \item \terme{Perspective} compatible.
    \item \terme{Échelle} cohérente.
    \item \terme{Style} compatible.
    \item \terme{Qualité} comparable.
\end{itemize}

\subsubsection{La pertinence}

Les photographies doivent être \terme{pertinentes} par rapport au concept :

\begin{itemize}
    \item \terme{Correspondre} à l'idée.
    \item \terme{Correspondre} au message.
    \item \terme{Correspondre} au ton.
    \item \terme{Correspondre} à l'univers visuel.
\end{itemize}

\subsubsection{La qualité}

Les photographies doivent être de \terme{qualité} :

\begin{itemize}
    \item \terme{Nettes} : pas de flou.
    \item \terme{Bien} exposées.
    \item \terme{Bien} contrastées.
    \item \terme{Haute} résolution.
    \item \terme{Sans} artefacts.
\end{itemize}

\subsubsection{Les droits}

Les photographies doivent être \terme{libres de droits} :

\begin{itemize}
    \item \terme{Vérifier} la licence.
    \item \terme{Respecter} les conditions d'utilisation.
    \item \terme{Citer} les auteurs si nécessaire.
    \item \terme{Éviter} les images protégées.
\end{itemize}

\subsection{Sélectionner les photographies}

\subsubsection{Étape 1 : Définir les besoins}

\begin{itemize}
    \item \terme{Lister} les éléments nécessaires (produit, environnement, personnes, textures).
    \item \terme{Définir} les critères (résolution, lumière, perspective, style).
    \item \terme{Définir} le nombre d'images.
\end{itemize}

\subsubsection{Étape 2 : Rechercher}

\begin{itemize}
    \item \terme{Chercher} sur les banques d'images.
    \item \terme{Chercher} des images personnelles.
    \item \terme{Chercher} des images générées.
    \item \terme{Utiliser} des mots-clés précis.
\end{itemize}

\subsubsection{Étape 3 : Sélectionner}

\begin{itemize}
    \item \terme{Sélectionner} les images les plus pertinentes.
    \item \terme{Vérifier} la compatibilité.
    \item \terme{Vérifier} la qualité.
    \item \terme{Vérifier} les droits.
\end{itemize}

\subsubsection{Étape 4 : Télécharger}

\begin{itemize}
    \item \terme{Télécharger} les images en haute résolution.
    \item \terme{Organiser} les fichiers.
    \item \terme{Nommer} les fichiers.
\end{itemize}

\subsubsection{Étape 5 : Vérifier}

\begin{itemize}
    \item \terme{Vérifier} la résolution.
    \item \terme{Vérifier} la netteté.
    \item \terme{Vérifier} la cohérence.
\end{itemize}

\subsection{Les bonnes pratiques}

\begin{itemize}
    \item \terme{Définir} les besoins.
    \item \terme{Rechercher} sur plusieurs sources.
    \item \terme{Sélectionner} avec soin.
    \item \terme{Vérifier} la compatibilité.
    \item \terme{Vérifier} la qualité.
    \item \terme{Vérifier} les droits.
    \item \terme{Organiser} les fichiers.
    \item \terme{Documenter} les sources.
\end{itemize}

\subsection*{Conclusion de la section}

La \terme{recherche de photographies} est essentielle pour constituer l'univers visuel. Les sources sont : banques d'images, photographies personnelles, images générées. Les critères sont : compatibilité, pertinence, qualité, droits. La sélection suit plusieurs étapes : définir les besoins, rechercher, sélectionner, télécharger, vérifier. Il est essentiel de définir les besoins, de vérifier la compatibilité, de vérifier les droits.

Maintenant que nous savons rechercher et sélectionner les photographies, nous allons, dans la section suivante, apprendre à \terme{construire la composition}.

% ============================================================
\section{Construction de la composition}
% ============================================================

\subsection{Pourquoi la composition ?}

La \terme{composition} est essentielle pour :

\begin{itemize}
    \item \terme{Structurer} le visuel.
    \item \terme{Transmettre} le message.
    \item \terme{Créer} une hiérarchie.
    \item \terme{Assurer} la lisibilité.
    \item \terme{Produire} un impact.
    \item \terme{Respecter} le concept.
\end{itemize}

\subsection{Les principes de composition}

\subsubsection{La hiérarchie}

La \terme{hiérarchie} organise les éléments : produit principal, éléments secondaires, environnement, typographie.

\subsubsection{L'équilibre}

L'\terme{équilibre} répartit les masses visuelles : symétrique ou asymétrique.

\subsubsection{Le contraste}

Le \terme{contraste} distingue les éléments : taille, couleur, valeur, position.

\subsubsection{La profondeur}

La \terme{profondeur} crée l'espace : plans, perspective, flou, atmosphère.

\subsubsection{Le point focal}

Le \terme{point focal} est l'élément principal : le produit, le visage, le message.

\subsubsection{L'espace négatif}

L'\terme{espace négatif} aère la composition : il met en valeur les éléments.

\subsection{Construire la composition}

\subsubsection{Étape 1 : Définir le point focal}

\begin{itemize}
    \item \terme{Définir} l'élément principal.
    \item \terme{Définir} sa position.
    \item \terme{Définir} sa taille.
\end{itemize}

\subsubsection{Étape 2 : Placer les éléments}

\begin{itemize}
    \item \terme{Placer} le produit au point focal.
    \item \terme{Placer} l'environnement.
    \item \terme{Placer} les éléments secondaires.
    \item \terme{Placer} la typographie.
\end{itemize}

\subsubsection{Étape 3 : Organiser les plans}

\begin{itemize}
    \item \terme{Définir} l'avant-plan.
    \item \terme{Définir} le plan moyen.
    \item \terme{Définir} l'arrière-plan.
    \item \terme{Créer} la profondeur.
\end{itemize}

\subsubsection{Étape 4 : Équilibrer}

\begin{itemize}
    \item \terme{Vérifier} l'équilibre.
    \item \terme{Ajuster} les positions.
    \item \terme{Ajuster} les tailles.
\end{itemize}

\subsubsection{Étape 5 : Vérifier}

\begin{itemize}
    \item \terme{Vérifier} la lisibilité.
    \item \terme{Vérifier} l'impact.
    \item \terme{Vérifier} la cohérence.
\end{itemize}

\subsection{Les bonnes pratiques}

\begin{itemize}
    \item \terme{Définir} le point focal.
    \item \terme{Placer} les éléments.
    \item \terme{Organiser} les plans.
    \item \terme{Équilibrer} la composition.
    \item \terme{Respecter} la hiérarchie.
    \item \terme{Utiliser} l'espace négatif.
    \item \terme{Vérifier} la lisibilité.
    \item \terme{Vérifier} l'impact.
\end{itemize}

\subsection*{Conclusion de la section}

La \terme{composition} est essentielle pour structurer le visuel, transmettre le message, créer une hiérarchie, assurer la lisibilité. Elle repose sur la hiérarchie, l'équilibre, le contraste, la profondeur, le point focal, l'espace négatif. La construction suit plusieurs étapes : définir le point focal, placer les éléments, organiser les plans, équilibrer, vérifier. Il est essentiel de définir le point focal, de placer les éléments, d'organiser les plans, d'équilibrer.

Maintenant que nous savons construire la composition, nous allons, dans la section suivante, apprendre à \terme{réaliser un détourage avancé}.

% ============================================================
\section{Détourage avancé}
% ============================================================

\subsection{Pourquoi un détourage avancé ?}

Un \terme{détourage avancé} est essentiel pour :

\begin{itemize}
    \item \terme{Isoler} le produit ou le sujet.
    \item \terme{Intégrer} le produit dans la composition.
    \item \terme{Supprimer} les éléments indésirables.
    \item \terme{Assurer} un résultat propre.
    \item \terme{Préserver} les détails.
    \item \terme{Éviter} les halos.
\end{itemize}

Un détourage approximatif se voit immédiatement et ruine le photomontage.

\subsection{Les outils de détourage avancé}

\subsubsection{L'outil Sélection d'objet}

\begin{itemize}
    \item \terme{Sélectionner} l'outil Sélection d'objet.
    \item \terme{Cliquer-glisser} autour du sujet.
    \item \terme{Photoshop} détecte le sujet.
    \item \terme{Affiner} si nécessaire.
\end{itemize}

\subsubsection{L'outil Sélection rapide}

\begin{itemize}
    \item \terme{Sélectionner} l'outil Sélection rapide.
    \item \terme{Peindre} sur le sujet.
    \item \terme{La sélection} s'étend automatiquement.
    \item \terme{Corriger} avec \texttt{Alt}.
\end{itemize}

\subsubsection{L'outil Plume}

\begin{itemize}
    \item \terme{Sélectionner} l'outil Plume.
    \item \terme{Tracer} le contour du sujet.
    \item \terme{Convertir} le tracé en sélection.
    \item \terme{Affiner}.
\end{itemize}

\subsubsection{Sélection > Sujet}

\begin{itemize}
    \item \terme{Utiliser} \texttt{Sélection > Sujet}.
    \item \terme{Photoshop} détecte le sujet principal.
    \item \terme{Affiner} si nécessaire.
\end{itemize}

\subsubsection{Sélectionner et masquer}

\begin{itemize}
    \item \terme{Ouvrir} \texttt{Sélection > Sélectionner et masquer}.
    \item \terme{Ajuster} le contour progressif.
    \item \terme{Ajuster} le lissage.
    \item \terme{Ajuster} la contraction/dilatation.
    \item \terme{Utiliser} \emph{Détecter les cheveux}.
    \item \terme{Sortir} vers un masque de fusion.
\end{itemize}

\subsection{Réaliser un détourage avancé}

\subsubsection{Étape 1 : Analyser le sujet}

\begin{itemize}
    \item \terme{Contours} nets ou flous.
    \item \terme{Cheveux, poils, textures}.
    \item \terme{Transparences, reflets}.
    \item \terme{Contraste} avec le fond.
\end{itemize}

\subsubsection{Étape 2 : Choisir la méthode}

\begin{itemize}
    \item \terme{Fond uni} : Baguette magique.
    \item \terme{Contours nets} : Plume, Lasso polygonal.
    \item \terme{Contours complexes} : Sélection d'objet, Sélection rapide.
    \item \terme{Portraits} : Sélection > Sujet + Sélectionner et masquer.
\end{itemize}

\subsubsection{Étape 3 : Créer la sélection}

\begin{itemize}
    \item \terme{Utiliser} l'outil choisi.
    \item \terme{Affiner} avec \emph{Sélectionner et masquer}.
\end{itemize}

\subsubsection{Étape 4 : Créer le masque}

\begin{itemize}
    \item \terme{Sortir} vers un masque de fusion.
    \item \terme{Vérifier} le résultat.
    \item \terme{Corriger} avec le pinceau.
\end{itemize}

\subsubsection{Étape 5 : Vérifier}

\begin{itemize}
    \item \terme{Zoomer} pour vérifier les contours.
    \item \terme{Vérifier} sur fond contrasté.
    \item \terme{Vérifier} les halos.
    \item \terme{Corriger} si nécessaire.
\end{itemize}

\subsection{Les techniques avancées}

\subsubsection{Le détourage des cheveux}

\begin{itemize}
    \item \terme{Utiliser} \emph{Détecter les cheveux} dans \emph{Sélectionner et masquer}.
    \item \terme{Ajuster} le rayon.
    \item \terme{Corriger} manuellement.
\end{itemize}

\subsubsection{Le détourage des objets transparents}

\begin{itemize}
    \item \terme{Utiliser} des masques de fusion avec dégradés.
    \item \terme{Combiner} avec des modes de fusion.
    \item \terme{Travailler} les reflets.
\end{itemize}

\subsubsection{Le détourage des objets flous}

\begin{itemize}
    \item \terme{Utiliser} un contour progressif élevé.
    \item \terme{Adoucir} les bords.
    \item \terme{Utiliser} un pinceau doux.
\end{itemize}

\subsection{Les bonnes pratiques}

\begin{itemize}
    \item \terme{Analyser} le sujet.
    \item \terme{Choisir} la bonne méthode.
    \item \terme{Affiner} avec \emph{Sélectionner et masquer}.
    \item \terme{Travailler} sur un masque.
    \item \terme{Corriger} avec le pinceau.
    \item \terme{Vérifier} sur fond contrasté.
    \item \terme{Zoomer} pour les détails.
    \item \terme{Préserver} la qualité.
\end{itemize}

\subsection*{Conclusion de la section}

Le \terme{détourage avancé} est essentiel pour isoler le produit ou le sujet. Les outils sont : Sélection d'objet, Sélection rapide, Plume, Sélection > Sujet, Sélectionner et masquer. La méthode suit plusieurs étapes : analyser le sujet, choisir la méthode, créer la sélection, créer le masque, vérifier. Il est essentiel d'analyser, de choisir la bonne méthode, d'affiner, de vérifier.

Maintenant que nous savons réaliser un détourage avancé, nous allons, dans la section suivante, apprendre à \terme{harmoniser la lumière}.

% ============================================================
\section{Harmonisation de la lumière}
% ============================================================

\subsection{Pourquoi harmoniser la lumière ?}

L'\terme{harmonisation de la lumière} est essentielle pour :

\begin{itemize}
    \item \terme{Intégrer} le produit dans l'environnement.
    \item \terme{Créer} une cohérence visuelle.
    \item \terme{Donner} du réalisme.
    \item \terme{Éviter} les incohérences.
    \item \terme{Assurer} la crédibilité.
    \item \terme{Produire} un impact.
\end{itemize}

Une lumière incohérente trahit le montage.

\subsection{Les paramètres de la lumière}

\subsubsection{La direction}

La \terme{direction} de la lumière doit être cohérente : si la lumière vient de la gauche dans l'environnement, elle doit venir de la gauche sur le produit.

\subsubsection{L'intensité}

L'\terme{intensité} de la lumière doit être cohérente : si l'environnement est lumineux, le produit doit être lumineux.

\subsubsection{La couleur}

La \terme{couleur} de la lumière doit être cohérente : si l'environnement est chaud, le produit doit être chaud.

\subsubsection{La dureté}

La \terme{dureté} de la lumière doit être cohérente : si l'environnement a une lumière douce, le produit doit avoir une lumière douce.

\subsection{Harmoniser la lumière}

\subsubsection{Étape 1 : Analyser la lumière}

\begin{itemize}
    \item \terme{Identifier} la direction.
    \item \terme{Identifier} l'intensité.
    \item \terme{Identifier} la couleur.
    \item \terme{Identifier} la dureté.
\end{itemize}

\subsubsection{Étape 2 : Ajuster le produit}

\begin{itemize}
    \item \terme{Éclaircir} ou \terme{assombrir} le produit.
    \item \terme{Réchauffer} ou \terme{refroidir} le produit.
    \item \terme{Ajuster} le contraste.
\end{itemize}

\subsubsection{Étape 3 : Ajuster l'environnement}

\begin{itemize}
    \item \terme{Ajuster} la lumière de l'environnement.
    \item \terme{Ajouter} des lumières.
    \item \terme{Ajouter} des ombres.
\end{itemize}

\subsubsection{Étape 4 : Vérifier la cohérence}

\begin{itemize}
    \item \terme{Vérifier} que la lumière est cohérente.
    \item \terme{Vérifier} que le produit s'intègre.
    \item \terme{Vérifier} que le résultat est crédible.
\end{itemize}

\subsection{Les techniques d'harmonisation}

\subsubsection{Les calques de réglage}

\begin{itemize}
    \item \terme{Calque de réglage} \emph{Courbes}.
    \item \terme{Calque de réglage} \emph{Niveaux}.
    \item \terme{Calque de réglage} \emph{Balance des couleurs}.
    \item \terme{Calque de réglage} \emph{Teinte/Saturation}.
\end{itemize}

\subsubsection{Les masques}

\begin{itemize}
    \item \terme{Masquer} les réglages sur les zones non concernées.
    \item \terme{Peindre} en blanc sur les zones à ajuster.
\end{itemize}

\subsubsection{Les modes de fusion}

\begin{itemize}
    \item \terme{Multiplier} : assombrit.
    \item \terme{Écran} : éclaircit.
    \item \terme{Superposition} : augmente le contraste.
    \item \terme{Couleur} : applique une couleur.
\end{itemize}

\subsection{Les bonnes pratiques}

\begin{itemize}
    \item \terme{Analyser} la lumière.
    \item \terme{Ajuster} le produit.
    \item \terme{Ajuster} l'environnement.
    \item \terme{Utiliser} des calques de réglage.
    \item \terme{Masquer} les réglages.
    \item \terme{Vérifier} la cohérence.
    \item \terme{Vérifier} la crédibilité.
    \item \terme{Tester} sur différents écrans.
\end{itemize}

\subsection*{Conclusion de la section}

L'\terme{harmonisation de la lumière} est essentielle pour intégrer le produit dans l'environnement. Elle repose sur la direction, l'intensité, la couleur, la dureté. La méthode suit plusieurs étapes : analyser la lumière, ajuster le produit, ajuster l'environnement, vérifier la cohérence. Il est essentiel d'analyser, d'ajuster, de masquer, de vérifier.

Maintenant que nous savons harmoniser la lumière, nous allons, dans la section suivante, apprendre à \terme{gérer les ombres et les reflets}.

% ============================================================
\section{Gestion des ombres et des reflets}
% ============================================================

\subsection{Pourquoi gérer les ombres et les reflets ?}

La \terme{gestion des ombres et des reflets} est essentielle pour :

\begin{itemize}
    \item \terme{Ancrer} le produit dans l'environnement.
    \item \terme{Créer} du réalisme.
    \item \terme{Donner} de la profondeur.
    \item \terme{Éviter} l'effet de flottement.
    \item \terme{Assurer} la crédibilité.
    \item \terme{Produire} un impact.
\end{itemize}

Un produit sans ombre flotte ; un produit sans reflet paraît faux.

\subsection{Les types d'ombres}

\subsubsection{L'ombre portée}

L'\terme{ombre portée} est projetée par le produit sur une surface. Elle dépend de la direction, de l'intensité, de la couleur de la lumière.

\subsubsection{L'ombre propre}

L'\terme{ombre propre} est la zone du produit non éclairée. Elle donne du volume.

\subsubsection{L'ombre de contact}

L'\terme{ombre de contact} est l'ombre très sombre au contact du produit et de la surface. Elle ancre le produit.

\subsection{Les types de reflets}

\subsubsection{Le reflet direct}

Le \terme{reflet direct} est le reflet du produit sur une surface réfléchissante (verre, eau, métal).

\subsubsection{Le reflet diffus}

Le \terme{reflet diffus} est la lumière réfléchie par le produit sur l'environnement.

\subsubsection{Le reflet spéculaire}

Le \terme{reflet spéculaire} est la lumière réfléchie par une surface brillante.

\subsection{Gérer les ombres}

\subsubsection{Étape 1 : Analyser la lumière}

\begin{itemize}
    \item \terme{Identifier} la direction.
    \item \terme{Identifier} l'intensité.
    \item \terme{Identifier} la couleur.
    \item \terme{Identifier} la dureté.
\end{itemize}

\subsubsection{Étape 2 : Créer l'ombre portée}

\begin{itemize}
    \item \terme{Dupliquer} le produit.
    \item \terme{Remplir} de noir.
    \item \terme{Appliquer} un flou.
    \item \terme{Déformer} pour la perspective.
    \item \terme{Réduire} l'opacité.
    \item \terme{Appliquer} un mode de fusion (Multiplier).
\end{itemize}

\subsubsection{Étape 3 : Créer l'ombre propre}

\begin{itemize}
    \item \terme{Peindre} en noir (ou en couleur sombre) sur les zones non éclairées.
    \item \terme{Utiliser} un mode de fusion (Multiplier).
    \item \terme{Ajuster} l'opacité.
\end{itemize}

\subsubsection{Étape 4 : Créer l'ombre de contact}

\begin{itemize}
    \item \terme{Peindre} une zone sombre au contact.
    \item \terme{Appliquer} un léger flou.
    \item \terme{Utiliser} un mode de fusion (Multiplier).
\end{itemize}

\subsection{Gérer les reflets}

\subsubsection{Étape 1 : Créer le reflet direct}

\begin{itemize}
    \item \terme{Dupliquer} le produit.
    \item \terme{Retourner} verticalement.
    \item \terme{Déformer} pour la perspective.
    \item \terme{Appliquer} un flou.
    \item \terme{Réduire} l'opacité.
    \item \terme{Masquer} pour ne garder que le reflet.
\end{itemize}

\subsubsection{Étape 2 : Créer le reflet diffus}

\begin{itemize}
    \item \terme{Peindre} avec une couleur claire.
    \item \terme{Appliquer} un flou.
    \item \terme{Utiliser} un mode de fusion (Écran).
\end{itemize}

\subsubsection{Étape 3 : Créer le reflet spéculaire}

\begin{itemize}
    \item \terme{Peindre} une zone claire.
    \item \terme{Utiliser} un pinceau doux.
    \item \terme{Utiliser} un mode de fusion (Écran).
\end{itemize}

\subsection{Les bonnes pratiques}

\begin{itemize}
    \item \terme{Analyser} la lumière.
    \item \terme{Créer} l'ombre portée.
    \item \terme{Créer} l'ombre propre.
    \item \terme{Créer} l'ombre de contact.
    \item \terme{Créer} le reflet direct.
    \item \terme{Créer} le reflet diffus.
    \item \terme{Vérifier} la cohérence.
    \item \terme{Vérifier} la crédibilité.
\end{itemize}

\subsection*{Conclusion de la section}

La \terme{gestion des ombres et des reflets} est essentielle pour ancrer le produit, créer du réalisme, donner de la profondeur. Les ombres sont : portée, propre, contact. Les reflets sont : direct, diffus, spéculaire. La méthode suit plusieurs étapes : analyser la lumière, créer les ombres, créer les reflets, vérifier. Il est essentiel d'analyser, de créer les ombres, de créer les reflets, de vérifier.

Maintenant que nous savons gérer les ombres et les reflets, nous allons, dans la section suivante, apprendre à \terme{intégrer le produit}.

% ============================================================
\section{Intégration du produit}
% ============================================================

\subsection{Pourquoi intégrer le produit ?}

L'\terme{intégration du produit} est essentielle pour :

\begin{itemize}
    \item \terme{Créer} un visuel cohérent.
    \item \terme{Mettre} en valeur le produit.
    \item \terme{Transmettre} le message.
    \item \terme{Assurer} la crédibilité.
    \item \terme{Produire} un impact.
    \item \terme{Respecter} le concept.
\end{itemize}

Un produit mal intégré ruine le photomontage, même si le reste est réussi.

\subsection{Les étapes d'intégration}

\subsubsection{Étape 1 : Positionner le produit}

\begin{itemize}
    \item \terme{Placer} le produit au point focal.
    \item \terme{Ajuster} la taille.
    \item \terme{Ajuster} la position.
    \item \terme{Vérifier} la perspective.
\end{itemize}

\subsubsection{Étape 2 : Ajuster la lumière}

\begin{itemize}
    \item \terme{Harmoniser} la direction.
    \item \terme{Harmoniser} l'intensité.
    \item \terme{Harmoniser} la couleur.
    \item \terme{Harmoniser} la dureté.
\end{itemize}

\subsubsection{Étape 3 : Créer les ombres}

\begin{itemize}
    \item \terme{Créer} l'ombre portée.
    \item \terme{Créer} l'ombre propre.
    \item \terme{Créer} l'ombre de contact.
\end{itemize}

\subsubsection{Étape 4 : Créer les reflets}

\begin{itemize}
    \item \terme{Créer} le reflet direct.
    \item \terme{Créer} le reflet diffus.
    \item \terme{Créer} le reflet spéculaire.
\end{itemize}

\subsubsection{Étape 5 : Vérifier l'intégration}

\begin{itemize}
    \item \terme{Vérifier} la cohérence.
    \item \terme{Vérifier} la crédibilité.
    \item \terme{Vérifier} l'impact.
    \item \terme{Zoomer} pour les détails.
\end{itemize}

\subsection{Les techniques avancées}

\subsubsection{Le placement dans l'environnement}

\begin{itemize}
    \item \terme{Placer} le produit sur une surface.
    \item \terme{Placer} le produit devant un fond.
    \item \terme{Placer} le produit dans un contexte.
\end{itemize}

\subsubsection{L'ajout d'éléments}

\begin{itemize}
    \item \terme{Ajouter} des éléments autour du produit.
    \item \terme{Ajouter} des textures.
    \item \terme{Ajouter} des effets.
\end{itemize}

\subsubsection{Le rendu final}

\begin{itemize}
    \item \terme{Ajuster} le contraste.
    \item \terme{Ajuster} la saturation.
    \item \terme{Ajuster} la netteté.
    \item \terme{Ajouter} un vignettage.
\end{itemize}

\subsection{Les bonnes pratiques}

\begin{itemize}
    \item \terme{Positionner} le produit.
    \item \terme{Ajuster} la lumière.
    \item \terme{Créer} les ombres.
    \item \terme{Créer} les reflets.
    \item \terme{Vérifier} l'intégration.
    \item \terme{Ajouter} des éléments.
    \item \terme{Ajuster} le rendu final.
    \item \terme{Vérifier} la cohérence.
\end{itemize}

\subsection*{Conclusion de la section}

L'\terme{intégration du produit} est essentielle pour créer un visuel cohérent. Elle repose sur le positionnement, la lumière, les ombres, les reflets. La méthode suit plusieurs étapes : positionner le produit, ajuster la lumière, créer les ombres, créer les reflets, vérifier. Il est essentiel de positionner, d'ajuster la lumière, de créer les ombres et les reflets, de vérifier.

Maintenant que nous savons intégrer le produit, nous allons, dans la section suivante, apprendre à \terme{traiter la colorimétrie}.

% ============================================================
\section{Traitement colorimétrique}
% ============================================================

\subsection{Pourquoi le traitement colorimétrique ?}

Le \terme{traitement colorimétrique} est essentiel pour :

\begin{itemize}
    \item \terme{Unifier} les couleurs.
    \item \terme{Créer} une ambiance.
    \item \terme{Transmettre} une émotion.
    \item \terme{Respecter} le concept.
    \item \terme{Assurer} la cohérence.
    \item \terme{Produire} un impact.
\end{itemize}

Un traitement colorimétrique réussi donne une identité forte au photomontage.

\subsection{Les paramètres colorimétriques}

\subsubsection{La teinte}

La \terme{teinte} définit la couleur dominante. Elle peut être :

\begin{itemize}
    \item \terme{Chaude} : rouges, oranges, jaunes.
    \item \terme{Froide} : bleus, verts, violets.
    \item \terme{Neutre} : gris, blancs.
\end{itemize}

\subsubsection{La saturation}

La \terme{saturation} définit l'intensité des couleurs. Elle peut être :

\begin{itemize}
    \item \terme{Vive} : couleurs intenses.
    \item \terme{Douce} : couleurs désaturées.
    \item \terme{Monochrome} : une seule couleur.
\end{itemize}

\subsubsection{La luminosité}

La \terme{luminosité} définit la clarté des couleurs. Elle peut être :

\begin{itemize}
    \item \terme{Claire} : tons clairs.
    \item \terme{Sombre} : tons sombres.
    \item \terme{Contrastée} : forts contrastes.
\end{itemize}

\subsection{Traiter la colorimétrie}

\subsubsection{Étape 1 : Définir l'ambiance}

\begin{itemize}
    \item \terme{Définir} l'ambiance (chaleureuse, froide, neutre).
    \item \terme{Définir} l'émotion (joie, calme, tension).
    \item \terme{Définir} le style (moderne, rétro, naturel).
\end{itemize}

\subsubsection{Étape 2 : Appliquer les réglages}

\begin{itemize}
    \item \terme{Calque de réglage} \emph{Balance des couleurs}.
    \item \terme{Calque de réglage} \emph{Teinte/Saturation}.
    \item \terme{Calque de réglage} \emph{Courbes}.
    \item \terme{Calque de réglage} \emph{Vibrance}.
\end{itemize}

\subsubsection{Étape 3 : Harmoniser}

\begin{itemize}
    \item \terme{Appliquer} les réglages à l'ensemble du visuel.
    \item \terme{Masquer} les réglages sur les zones non concernées.
    \item \terme{Vérifier} la cohérence.
\end{itemize}

\subsubsection{Étape 4 : Vérifier}

\begin{itemize}
    \item \terme{Vérifier} la cohérence.
    \item \terme{Vérifier} l'ambiance.
    \item \terme{Vérifier} l'émotion.
    \item \terme{Vérifier} la crédibilité.
\end{itemize}

\subsection{Les techniques avancées}

\subsubsection{Le split toning}

Le \terme{split toning} applique des couleurs différentes aux ombres et aux hautes lumières. Il crée une ambiance riche.

\subsubsection{Le color grading}

Le \terme{color grading} applique une palette de couleurs spécifique pour créer une ambiance cinématographique.

\subsubsection{Le virage partiel}

Le \terme{virage partiel} applique des couleurs différentes dans les ombres, les tons moyens, les hautes lumières.

\subsection{Les bonnes pratiques}

\begin{itemize}
    \item \terme{Définir} l'ambiance.
    \item \terme{Appliquer} les réglages.
    \item \terme{Harmoniser} l'ensemble.
    \item \terme{Masquer} les réglages.
    \item \terme{Vérifier} la cohérence.
    \item \terme{Vérifier} l'ambiance.
    \item \terme{Vérifier} l'émotion.
    \item \terme{Tester} sur différents écrans.
\end{itemize}

\subsection*{Conclusion de la section}

Le \terme{traitement colorimétrique} est essentiel pour unifier les couleurs, créer une ambiance, transmettre une émotion. Il repose sur la teinte, la saturation, la luminosité. La méthode suit plusieurs étapes : définir l'ambiance, appliquer les réglages, harmoniser, vérifier. Il est essentiel de définir l'ambiance, d'appliquer les réglages, de vérifier la cohérence.

Maintenant que nous savons traiter la colorimétrie, nous allons, dans la section suivante, apprendre à \terme{intégrer la typographie}.

% ============================================================
\section{Intégration de la typographie}
% ============================================================

\subsection{Pourquoi intégrer la typographie ?}

L'\terme{intégration de la typographie} est essentielle pour :

\begin{itemize}
    \item \terme{Transmettre} le message.
    \item \terme{Identifier} la marque.
    \item \terme{Compléter} le visuel.
    \item \terme{Créer} une hiérarchie.
    \item \terme{Assurer} la cohérence.
    \item \terme{Produire} un impact.
\end{itemize}

Une typographie mal intégrée ruine le photomontage, même si le visuel est réussi.

\subsection{Les éléments typographiques}

\subsubsection{Le titre}

Le \terme{titre} est le message principal. Il doit être :

\begin{itemize}
    \item \terme{Court} : 3 à 7 mots.
    \item \terme{Percutant} : fort, mémorable.
    \item \terme{Lisible} : grand, contrasté.
    \item \terme{Cohérent} : avec l'ambiance.
\end{itemize}

\subsubsection{Le sous-titre}

Le \terme{sous-titre} complète le titre. Il doit être :

\begin{itemize}
    \item \terme{Clair} : compréhensible.
    \item \terme{Secondaire} : moins visible que le titre.
    \item \terme{Cohérent} : avec le titre.
\end{itemize}

\subsubsection{Le texte}

Le \terme{texte} donne les informations détaillées. Il doit être :

\begin{itemize}
    \item \terme{Lisible} : taille suffisante.
    \item \terme{Structuré} : organisé.
    \item \terme{Cohérent} : avec le reste.
\end{itemize}

\subsubsection{Le logo}

Le \terme{logo} identifie la marque. Il doit être :

\begin{itemize}
    \item \terme{Visible} : mais discret.
    \item \terme{Cohérent} : avec la charte.
    \item \terme{Placé} judicieusement.
\end{itemize}

\subsection{Intégrer la typographie}

\subsubsection{Étape 1 : Choisir la typographie}

\begin{itemize}
    \item \terme{Choisir} une police adaptée à l'ambiance.
    \item \terme{Choisir} une police lisible.
    \item \terme{Choisir} une police cohérente.
\end{itemize}

\subsubsection{Étape 2 : Placer le texte}

\begin{itemize}
    \item \terme{Placer} le titre au point focal.
    \item \terme{Placer} le sous-titre.
    \item \terme{Placer} le texte.
    \item \terme{Placer} le logo.
\end{itemize}

\subsubsection{Étape 3 : Appliquer les effets}

\begin{itemize}
    \item \terme{Appliquer} une ombre portée (si nécessaire).
    \item \terme{Appliquer} un contour (si nécessaire).
    \item \terme{Appliquer} un mode de fusion (si nécessaire).
\end{itemize}

\subsubsection{Étape 4 : Vérifier}

\begin{itemize}
    \item \terme{Vérifier} la lisibilité.
    \item \terme{Vérifier} la hiérarchie.
    \item \terme{Vérifier} la cohérence.
    \item \terme{Vérifier} l'intégration.
\end{itemize}

\subsection{Les bonnes pratiques}

\begin{itemize}
    \item \terme{Choisir} une typographie adaptée.
    \item \terme{Placer} le texte judicieusement.
    \item \terme{Appliquer} les effets avec parcimonie.
    \item \terme{Vérifier} la lisibilité.
    \item \terme{Vérifier} la hiérarchie.
    \item \terme{Vérifier} la cohérence.
    \item \terme{Vérifier} l'intégration.
    \item \terme{Tester} sur différents supports.
\end{itemize}

\subsection*{Conclusion de la section}

L'\terme{intégration de la typographie} est essentielle pour transmettre le message, identifier la marque, compléter le visuel. Elle comprend le titre, le sous-titre, le texte, le logo. La méthode suit plusieurs étapes : choisir la typographie, placer le texte, appliquer les effets, vérifier. Il est essentiel de choisir une typographie adaptée, de placer le texte judicieusement, de vérifier la lisibilité.

Maintenant que nous savons intégrer la typographie, nous allons, dans la dernière section de ce chapitre, apprendre à \terme{finaliser et contrôler la qualité}.

% ============================================================
\section{Finition et contrôle qualité}
% ============================================================

\subsection{Pourquoi finaliser et contrôler la qualité ?}

La \terme{finalisation} et le \terme{contrôle qualité} sont essentiels pour :

\begin{itemize}
    \item \terme{Assurer} la qualité du résultat.
    \item \terme{Détecter} les erreurs.
    \item \terme{Corriger} les problèmes.
    \item \terme{Préparer} la livraison.
    \item \terme{Préserver} la réputation.
    \item \terme{Satisfaire} le client.
\end{itemize}

Un photomontage mal finalisé peut ruiner tout le travail.

\subsection{La finalisation}

\subsubsection{Étape 1 : Vérifier l'ensemble}

\begin{itemize}
    \item \terme{Vérifier} la composition.
    \item \terme{Vérifier} l'intégration.
    \item \terme{Vérifier} la cohérence.
    \item \terme{Vérifier} l'impact.
\end{itemize}

\subsubsection{Étape 2 : Ajuster les détails}

\begin{itemize}
    \item \terme{Ajuster} le contraste global.
    \item \terme{Ajuster} la saturation.
    \item \terme{Ajuster} la netteté.
    \item \terme{Ajouter} un vignettage léger.
\end{itemize}

\subsubsection{Étape 3 : Préparer l'export}

\begin{itemize}
    \item \terme{Choisir} le format (JPEG, PNG, PDF).
    \item \terme{Choisir} la résolution (72 ppi pour le web, 300 ppi pour l'impression).
    \item \terme{Choisir} le profil (sRGB pour le web, CMJN pour l'impression).
    \item \terme{Configurer} la compression.
\end{itemize}

\subsubsection{Étape 4 : Nommer et organiser}

\begin{itemize}
    \item \terme{Nommer} le fichier clairement.
    \item \terme{Organiser} dans des dossiers.
    \item \terme{Conserver} les fichiers sources.
\end{itemize}

\subsection{Le contrôle qualité}

\subsubsection{La lisibilité}

\begin{itemize}
    \item \terme{Vérifier} que le texte est lisible.
    \item \terme{Vérifier} le contraste.
    \item \terme{Vérifier} la taille.
\end{itemize}

\subsubsection{La cohérence}

\begin{itemize}
    \item \terme{Vérifier} la cohérence des couleurs.
    \item \terme{Vérifier} la cohérence de la lumière.
    \item \terme{Vérifier} la cohérence des ombres.
\end{itemize}

\subsubsection{La crédibilité}

\begin{itemize}
    \item \terme{Vérifier} que l'intégration est crédible.
    \item \terme{Vérifier} que les perspectives sont cohérentes.
    \item \terme{Vérifier} que les échelles sont correctes.
\end{itemize}

\subsubsection{La qualité technique}

\begin{itemize}
    \item \terme{Vérifier} la résolution.
    \item \terme{Vérifier} la netteté.
    \item \terme{Vérifier} les artefacts.
    \item \terme{Vérifier} le poids.
\end{itemize}

\subsection{Les bonnes pratiques}

\begin{itemize}
    \item \terme{Vérifier} l'ensemble.
    \item \terme{Ajuster} les détails.
    \item \terme{Préparer} l'export.
    \item \terme{Nommer} et organiser.
    \item \terme{Vérifier} la lisibilité.
    \item \terme{Vérifier} la cohérence.
    \item \terme{Vérifier} la crédibilité.
    \item \terme{Vérifier} la qualité technique.
\end{itemize}

\subsection*{Conclusion de la section}

La \terme{finalisation} et le \terme{contrôle qualité} sont essentiels pour assurer la qualité du résultat, détecter les erreurs, préparer la livraison. La finalisation suit plusieurs étapes : vérifier l'ensemble, ajuster les détails, préparer l'export, nommer et organiser. Le contrôle qualité vérifie la lisibilité, la cohérence, la crédibilité, la qualité technique. Il est essentiel de vérifier, d'ajuster, de préparer l'export, de contrôler la qualité.

% ============================================================
\section*{Conclusion du chapitre}
\addcontentsline{toc}{section}{Conclusion du chapitre}
% ============================================================

Nous voici au terme du chapitre 26. Nous avons appris à maîtriser le photomontage publicitaire avancé :

\begin{itemize}
    \item \terme{Analyse du concept publicitaire} : idée, message, bénéfice, ton, univers visuel.
    \item \terme{Recherche et sélection des photographies} : sources, critères (compatibilité, pertinence, qualité, droits), méthode.
    \item \terme{Construction de la composition} : hiérarchie, équilibre, contraste, profondeur, point focal, espace négatif.
    \item \terme{Détourage avancé} : outils (Sélection d'objet, Sélection rapide, Plume, Sélection > Sujet, Sélectionner et masquer), techniques avancées (cheveux, transparences, flou).
    \item \terme{Harmonisation de la lumière} : direction, intensité, couleur, dureté ; calques de réglage, masques, modes de fusion.
    \item \terme{Gestion des ombres et des reflets} : ombres portée, propre, contact ; reflets direct, diffus, spéculaire.
    \item \terme{Intégration du produit} : positionnement, lumière, ombres, reflets, rendu final.
    \item \terme{Traitement colorimétrique} : teinte, saturation, luminosité ; split toning, color grading, virage partiel.
    \item \terme{Intégration de la typographie} : titre, sous-titre, texte, logo ; placement, effets.
    \item \terme{Finition et contrôle qualité} : vérification, ajustement, export, nommage.
\end{itemize}

Le photomontage publicitaire est l'art de créer une image qui n'existe pas, en assemblant plusieurs photographies, illustrations, textures, pour construire un univers visuel cohérent, séduisant, mémorable. Il est utilisé dans la publicité, le cinéma, l'édition, le packaging, le web. Il permet de mettre en scène un produit dans un environnement imaginaire, de créer une situation surréaliste, de raconter une histoire, de susciter une émotion.

Le photomontage publicitaire exige une maîtrise avancée de tous les outils de Photoshop, une compréhension fine du concept publicitaire, une rigueur dans l'intégration du produit, une créativité dans la mise en scène. C'est l'un des domaines les plus créatifs et les plus techniques du design graphique. Un photomontage réussi est un photomontage où l'on ne voit pas le montage : les images semblent avoir été prises ensemble, dans le même lieu, à la même lumière.

Ce chapitre t'a donné les bases du photomontage publicitaire avancé. Mais la maîtrise vient avec la pratique. N'hésite pas à t'entraîner sur des projets variés : publicités, affiches, couvertures, packagings. Chaque projet est une occasion de progresser.

Ce chapitre clôt la quatrième partie de ce cours, consacrée au \terme{design publicitaire et à la communication visuelle}. Tu disposes désormais des compétences nécessaires pour créer des affiches, des flyers, des visuels pour les réseaux sociaux, des identités visuelles, des photomontages publicitaires. Ces compétences sont essentielles, mais elles ne suffisent pas : il faut aussi maîtriser le design éditorial, la préparation à l'impression, la gestion de projet, la tarification. C'est l'objet des parties suivantes.

Dans le chapitre suivant, nous aborderons les \terme{fondamentaux de la production imprimée}. Tu apprendras à comprendre les étapes de fabrication d'un imprimé, à distinguer RVB et CMJN, à gérer la résolution et les dimensions, à respecter les marges de sécurité, à préparer les fonds perdus et les traits de coupe, à gérer les profils colorimétriques, à incorporer les polices, à contrôler les images et les liens, à préparer les fichiers pour l'imprimeur. Ces compétences te permettront de préparer des fichiers professionnels pour l'impression.

% ============================================================
\section*{Projet de la Partie IV}
\addcontentsline{toc}{section}{Projet de la Partie IV}
% ============================================================

\subsection*{Réaliser une campagne publicitaire cohérente}

\subsubsection*{Objectif}

Mettre en pratique l'ensemble des compétences acquises dans la Partie IV : méthodologie, affiches, flyers, réseaux sociaux, branding, photomontage.

\subsubsection*{Sujet}

Réaliser une campagne publicitaire cohérente pour une entreprise fictive. Le sujet est libre, mais il doit être réaliste et cohérent. Exemples : lancement d'un produit, promotion d'un événement, campagne de sensibilisation, développement d'une marque.

\subsubsection*{Contraintes}

\begin{itemize}
    \item \terme{Concept} : définir un concept publicitaire fort.
    \item \terme{Affiche principale} : créer une affiche A3 (photographique, illustrative, ou typographique).
    \item \terme{Déclinaisons} : créer au moins 3 déclinaisons (flyer A5, post Instagram, story, bannière web).
    \item \terme{Photomontage} : intégrer au moins un photomontage publicitaire (produit, ambiance, situation).
    \item \terme{Identité} : respecter une identité visuelle cohérente (logo, couleurs, typographies).
    \item \terme{Charte} : documenter la campagne (concept, choix graphiques, applications).
    \item \terme{Présentation} : préparer une présentation de la campagne (démarche, visuels, applications).
\end{itemize}

\subsubsection*{Étapes}

\begin{enumerate}
    \item \terme{Brief} : définir la marque, le produit, le public, le concept.
    \item \terme{Recherche} : références, moodboard, esquisses.
    \item \terme{Création} : identité visuelle, affiche principale, photomontage.
    \item \terme{Déclinaison} : flyer, réseaux sociaux, bannière.
    \item \terme{Finalisation} : contrôle qualité, export.
    \item \terme{Présentation} : document de présentation, charte de campagne.
\end{enumerate}

\subsubsection*{Auto-évaluation}

\begin{table}[h!]
\centering
\begin{tabular}{@{}lccccc@{}}
\toprule
\textbf{Critère} & \textbf{1} & \textbf{2} & \textbf{3} & \textbf{4} & \textbf{5} \\
\midrule
Concept clair et fort & \(\square\) & \(\square\) & \(\square\) & \(\square\) & \(\square\) \\
Recherche pertinente & \(\square\) & \(\square\) & \(\square\) & \(\square\) & \(\square\) \\
Identité cohérente & \(\square\) & \(\square\) & \(\square\) & \(\square\) & \(\square\) \\
Affiche efficace & \(\square\) & \(\square\) & \(\square\) & \(\square\) & \(\square\) \\
Photomontage réussi & \(\square\) & \(\square\) & \(\square\) & \(\square\) & \(\square\) \\
Déclinaisons adaptées & \(\square\) & \(\square\) & \(\square\) & \(\square\) & \(\square\) \\
Cohérence de la campagne & \(\square\) & \(\square\) & \(\square\) & \(\square\) & \(\square\) \\
Lisibilité sur mobile & \(\square\) & \(\square\) & \(\square\) & \(\square\) & \(\square\) \\
Qualité technique & \(\square\) & \(\square\) & \(\square\) & \(\square\) & \(\square\) \\
Présentation claire & \(\square\) & \(\square\) & \(\square\) & \(\square\) & \(\square\) \\
Qualité globale & \(\square\) & \(\square\) & \(\square\) & \(\square\) & \(\square\) \\
\bottomrule
\end{tabular}
\caption{Grille d'auto-évaluation du projet de la Partie IV.}
\end{table}

\subsubsection*{Livrables}

\begin{itemize}
    \item \terme{Fichiers sources} : PSD, AI, avec calques organisés.
    \item \terme{Exports} : affiche (PDF, JPEG), flyer (PDF, JPEG), posts (PNG, JPEG), bannière (PNG).
    \item \terme{Charte de campagne} : document de présentation.
    \item \terme{Présentation} : document de présentation de la campagne.
\end{itemize}

\vspace{2em}

\begin{center}
\textit{Fin du chapitre 26 --- Photomontage publicitaire avancé.}

\vspace{1em}

\textit{Fin de la Partie IV --- Design publicitaire et communication visuelle.}

\vspace{1em}

\textbf{Prochain chapitre : Chapitre 27 --- Fondamentaux de la production imprimée.}
\end{center}